\documentclass[11pt,a4paper]{article}
\usepackage[margin=2.3cm]{geometry}
\usepackage{amsmath,amssymb,booktabs,graphicx,longtable,xcolor}
\usepackage[colorlinks=true,linkcolor=blue!45!black,urlcolor=blue!45!black]{hyperref}
\usepackage{titlesec}
\titleformat{\section}{\large\bfseries}{\thesection}{0.6em}{}
\setlength{\parskip}{0.35em}
\newcommand{\pip}{\mathrm{PIP}}
\title{\bf Fine-mapping across sparse and sparse-plus-infinitesimal architectures\\[3pt]
\large Iteration 004 benchmark --- complete analysis, all six strata}
\author{fmbenchmark}\date{\today}
\begin{document}\maketitle

\begin{abstract}\noindent
This report analyses the complete Iteration 004 fine-mapping benchmark: 45 design
rows, 11{,}250 scenarios, 112{,}500 region-fits per method, 18 methods, in-sample
LD throughout. It spans two genetic architectures --- a purely sparse model and a
sparse-plus-infinitesimal model --- crossed with three annotation conditions,
giving six strata that are never pooled. Every quantity derived from posterior
inclusion probabilities is defined explicitly before use, including its
aggregation level, its tie handling and the noise correction applied to the
variance decomposition. The principal findings are: (i) the number of causal
variants $S$ dominates the variance budget in every stratum; (ii) the
\emph{presence} of annotations matters and their \emph{strength} barely does;
(iii) the Functional BEATRICE ablation is monotone and highly significant in all
four annotated strata and \emph{exactly zero} in both unannotated ones, over
51{,}250 paired replicates; and (iv) the annotation-aware methods that lead under
the sparse model degrade sharply under misspecification, whereas the
region-sharing variant does not --- \texttt{fb\_xregion} rises from third to
second of eighteen and overtakes the entire PolyFun family.
\end{abstract}

\tableofcontents\bigskip

\section{Design}\label{sec:design}

\subsection{The two architectures}

Each simulated region contains exactly $S$ causal variants. The arms differ in
what else contributes to the phenotype.

\paragraph{Sparse.} Only the $S$ causal variants have effects. Effect sizes are
drawn $\beta\sim N(0,\sigma^2_\beta)$ and the residual variance is calibrated to
hit the target heritability exactly:
\[
\sigma^2_\varepsilon=\operatorname{Var}(X\beta)\,\frac{1-\phi}{\phi}
\quad\Longrightarrow\quad
\mathrm{PVE}=\frac{\operatorname{Var}(X\beta)}{\operatorname{Var}(X\beta)+\sigma^2_\varepsilon}=\phi .
\]

\paragraph{Sparse plus infinitesimal.} A second genetic component is added,
drawn on the \textbf{non-causal variants only}
($\alpha_{nc}\sim N(0,1/m_{nc})$, the BEATRICE formulation). The three variance
components are then scaled so that
\[
\operatorname{Var}(g_{\text{sparse}})=p_{\text{causal}}\cdot\phi,\qquad
\operatorname{Var}(g_{\text{inf}})=(1-p_{\text{causal}})\cdot\phi,\qquad
\sigma^2_\varepsilon=1-\phi .
\]
So $\phi$ is total heritability in both arms, and $p_{\text{causal}}$ is
\textbf{the fraction of heritability carried by the sparse component}. At
$p_{\text{causal}}=0.2$, four-fifths of the genetic signal is infinitesimal.

\paragraph{This has a consequence that governs how the second arm must be read.}
Under sparse-plus-infinitesimal, \emph{every non-causal variant has a genuine
non-zero effect}. The ground truth used for every metric in this report is the
set of $S$ sparse causal variants only. A method that correctly detects
infinitesimal signal is therefore scored as making a false discovery.
\textbf{All accuracy and calibration numbers in the second arm measure recovery
of the sparse component against a background of real but unlabelled signal},
which is the correct question for fine-mapping but is not the same question as
the first arm asks. Cross-arm comparisons of absolute FDR are meaningless for
that reason; comparisons of \emph{degradation} between arms are the informative
quantity, and are the subject of \S\ref{sec:misspec}.

\subsection{The grid}

\begin{center}\small
\begin{tabular}{llll}\toprule
& rows & annotation & enrichment folds\\\midrule
sparse & 1 & none & ---\\
& 4 & binary & 2.7, 5.4, 8.1, 10.8\\
& 4 & continuous & 2.7, 5.4, 8.1, 10.8\\
sparse + infinitesimal & 4 & none & --- \quad ($p_{\text{causal}}\in\{0.2,0.4,0.6,0.8\}$)\\
& 16 & binary & $4\times4$ enrichment $\times\ p_{\text{causal}}$\\
& 16 & continuous & $4\times4$\\\midrule
\multicolumn{4}{l}{45 rows $\times$ 250 scenarios $\times$ 10 regions $=$ 112{,}500 fits per method}\\
\bottomrule\end{tabular}\end{center}

Within every row: $S\in\{1,2,3,5,10\}$, $\phi\in\{0.05,0.1,0.2,0.4,0.6\}$,
10 iterations, 10 regions of nominal sizes
$\{500,500,750,750,1000,1000,1500,1500,2000,2000\}$, $n=5000$ individuals.

\paragraph{Causal variants are drawn with annotation-proportional probability.}
$w_j=\exp\!\big(A_j^\top\log\gamma\big)$, normalised, sampled without
replacement. \texttt{polyfun\_oracle} reconstructs exactly this $\pi_j$ from the
stored truth, which is what makes it a genuine ceiling rather than merely a
strong method.

\paragraph{Nominal region sizes are approximate at the top of the range.}
Regions are drawn from real 1000 Genomes loci, and many do not contain 2{,}000
usable variants. Measured actual counts:

\begin{center}\small
\begin{tabular}{lrrrr}\toprule
nominal & median actual & ratio & within 5\% of nominal\\\midrule
500  & 499  & 1.00 & 100\%\\
750  & 749  & 1.00 & 96\%\\
1000 & 999  & 1.00 & 97\%\\
1500 & 1483 & 0.99 & 62\%\\
2000 & \textbf{1590} & \textbf{0.80} & \textbf{26\%}\\
\bottomrule\end{tabular}\end{center}

The true swept range is therefore $499$--$1590$ ($3.2\times$), not
$500$--$2000$ ($4\times$). Since $\omega^2$ is range-dependent
(\S\ref{sec:sensev}), region size's already-small importance is measured over a
narrower span than its label suggests. The design remains valid: variation in
actual size \emph{within} a nominal class is differenced away by the
$\sigma^2_u$ estimator (\S\ref{sec:vhow}) and correctly subtracted. Only the
label is approximate.

\paragraph{The unit of analysis.} A \textbf{cell} is
$(\text{row},S,\phi,\text{nominal region size})$, aggregating
$10\ \text{iterations}\times2\ \text{regions of that size}=20$ fits. There are
$45\times125=5{,}625$ cells per method. Cell means feed every ANOVA;
individual fits feed the calibration counts; replicate-level rows feed the
paired contrasts.

\paragraph{Stratification is not optional.} Nothing is pooled across
architecture or annotation type. Beyond the obvious --- a method that cannot use
annotations is not comparable to one that can --- the unannotated arms serve as
negative controls, and the two architectures pose different questions
(\S\ref{sec:design}).
\section{From posterior inclusion probabilities to numbers}\label{sec:methods}

Every method returns, for each region, a vector
$\pip_1,\dots,\pip_p\in[0,1]$ --- one posterior inclusion probability per
variant --- together with the truth $y_j\in\{0,1\}$ known from the simulation.
All metrics below are functions of that pair. This section defines each one and
justifies the estimator, because several of the obvious implementations are
wrong in ways that silently change conclusions.

\subsection{Average precision (the ranking metric)}\label{sec:ap}

Average precision summarises the whole precision--recall curve without
committing to a threshold, which matters here because methods' PIP scales
differ by orders of magnitude: any fixed-threshold comparison partly measures
calibration rather than ranking.

Sort the variants by $\pip$ descending. Let $k$ index the distinct PIP values
(so that all variants sharing a value form one block), $\mathrm{tp}_k$ the
cumulative number of causals at the end of block $k$, $n_k$ the cumulative
number of variants, $\mathrm{prec}_k=\mathrm{tp}_k/n_k$ and
$\mathrm{rec}_k=\mathrm{tp}_k/S$. Then
\begin{equation}
\mathrm{AP}=\sum_k\bigl(\mathrm{rec}_k-\mathrm{rec}_{k-1}\bigr)\,\mathrm{prec}_k ,
\qquad \mathrm{rec}_0=0 .
\label{eq:ap}
\end{equation}

Three implementation points are load-bearing.

\textbf{It is a step integral, not a trapezoid.} Equation~\eqref{eq:ap}
evaluates precision at the operating point reached, matching
\texttt{sklearn.metrics.average\_precision\_score}. A trapezoidal rule
interpolates between operating points that the classifier never actually
occupies and is optimistically biased.

\textbf{Ties are resolved by block, not by index.} Several methods emit many
exactly-tied PIPs. Implementing AP as ``sort, then read off precision at the
causal's index'' lets the sort's tie-break --- here genomic position --- decide
the score. Because causal variants are drawn uniformly at random this would not
introduce a constant bias but would inject pure noise into precisely the
comparisons the analysis exists to make. Collapsing tied variants into one
block gives every member the same operating point, which is the only
tie-handling that is invariant to input order.

\textbf{AP is computed per fit and then averaged, never on a merged ranking.}
Pooling all variants from all regions into one ranking and computing a single
AP is a different and incorrect quantity: it rewards a method for getting the
\emph{relative calibration across regions} right, which is not what
region-level fine-mapping is asked to do. The two disagree materially. For two
fits with $\mathrm{AP}=1$ and $\mathrm{AP}=1/3$, the correct answer is the mean
$2/3\approx0.667$, whereas the merged ranking of the same data gives $0.700$.
Both values are checked against hand computation in the test suite.

A fit with no causal variants returns \texttt{NA} rather than $0$, and NAs are
excluded from means with the count reported alongside, so a method is never
credited or penalised for a fit that posed no question.

\subsection{Thresholded discovery: FDR and its honest companion}\label{sec:fdr}

For a threshold $t$ let $\mathcal{S}_t=\{j:\pip_j\ge t\}$, $n_t=|\mathcal{S}_t|$
and $\mathrm{tp}_t=\sum_{j\in\mathcal{S}_t}y_j$. Then
\begin{equation}
\mathrm{FDR}_t=\frac{n_t-\mathrm{tp}_t}{n_t},\qquad
\text{reported jointly with } n_t .
\end{equation}
Reporting FDR alone is not meaningful: a method that never crosses $t$ has
$\mathrm{FDR}_t=0$ by convention and $n_t=0$, and would otherwise win the
comparison by abstaining. The pair $(\mathrm{FDR}_t,n_t)$ is therefore always
presented together, and $n_t$ appears in the credible-set analysis
(\S\ref{sec:cs}) as the follow-up burden.

We also compute the \textbf{honesty ratio}. A Bayesian method that emits $\pip$
values is making a falsifiable claim about its own error rate: the expected
number of false positives among $\mathcal{S}_t$ is
$\widehat{\mathrm{FP}}_t=\sum_{j\in\mathcal{S}_t}(1-\pip_j)$. Comparing that to
the realised count gives
\begin{equation}
\rho_t=\frac{\mathrm{FP}_t}{\widehat{\mathrm{FP}}_t},
\end{equation}
so $\rho_t=1$ means the method's own uncertainty statement was accurate,
$\rho_t>1$ means it was overconfident by that factor. This is stricter than
FDR: it holds the method to \emph{its own} claim rather than to an external
threshold.

\subsection{Calibration}\label{sec:calibmethod}

Two complementary summaries are used, because each is separately gameable.

\textbf{Total mass ratio.} $\sum_j \pip_j / S$. A coherent posterior over which
variants are causal should place total mass equal to the expected number of
causals. A ratio of $2$ means the method claims twice as many causal variants
as exist.

\textbf{Banded reliability.} Variants are binned by PIP into
$[0,0.1)$, $[0.1,0.5)$, $[0.5,0.9)$, $[0.9,1]$. Within a bin we accumulate the
count $n_b$, the number truly causal $c_b$, and $\sum\pip$. Calibration
compares the realised rate $c_b/n_b$ to the claimed mean $\sum\pip/n_b$; this is
the reliability diagram in Figure~\ref{fig:rel}. We report
$\mathrm{ECE}_{hi}$, the expected calibration error restricted to bins with
$\pip\ge0.1$, because the $[0,0.1)$ bin contains the overwhelming majority of
variants and an unrestricted ECE is dominated by the trivially-correct
statement that almost nothing is causal.

\textbf{Why both, and why RES.} Consider a predictor emitting $\pip_j=S/p$ for
every variant. Its mass ratio is exactly $1$ and its calibration error is
exactly $0$ --- it is perfectly calibrated and completely useless. The Murphy
decomposition of the Brier score separates this out:
\begin{equation}
\mathrm{BS}=\underbrace{\mathrm{REL}}_{\text{miscalibration}}-\underbrace{\mathrm{RES}}_{\text{resolution}}+\underbrace{\mathrm{UNC}}_{\text{base rate}},
\end{equation}
where RES measures how far the method's conditional rates depart from the base
rate. The uniform predictor has $\mathrm{REL}=0$ \emph{and} $\mathrm{RES}=0$.
Resolution is therefore the component that detects calibration achieved by
saying nothing, and it is reported in Table~\ref{tab:main-sparse-continuous} and
Figure~\ref{fig:murphy} alongside REL.

\subsection{Credible sets}\label{sec:cs}

For a target mass $\alpha$, variants are added in decreasing PIP order until the
cumulative mass first reaches $\alpha$; the resulting set size is
$|CS_\alpha|$ and coverage is the indicator that the set contains a causal
variant, averaged over fits. When a posterior's \emph{total} mass never reaches
$\alpha$ the set is the whole region and a flag \texttt{set\_reached} records
it --- this is a result, not a missing value, and suppressing it would let a
maximally diffuse posterior masquerade as a confident small set. Size and
coverage are always read together (Figure~\ref{fig:cs}): a set can be made to
cover by being large, and small by abandoning coverage.

\subsection{The replication structure and the noise floor}\label{sec:noise}

The design is a split plot. Two variance components matter:
$\sigma^2_\varepsilon$, the between-iteration variance at fixed region, and
$\sigma^2_u$, the between-region-draw variance. A factor whose levels are
compared \emph{within} the same region draws (namely $S$ and $\phi$, which are
re-simulated per scenario) sees only $\sigma^2_\varepsilon$. A factor whose
levels necessarily use \emph{different} region draws --- region size, and any
term containing it --- additionally carries $\sigma^2_u$. Those are
``whole-plot'' terms.

For a term $T$ with $\mathrm{df}_T$ degrees of freedom, the expected sum of
squares under the null is $\mathrm{df}_T\lambda_T$ with
\begin{equation}
\lambda_T=\frac{\sigma^2_\varepsilon}{20}
+\mathbf{1}\{T\ \text{whole-plot}\}\cdot\frac{m\,\sigma^2_u}{2},
\qquad m=25 ,
\end{equation}
the $20$ being the fits per cell and the $2$ the regions per size class. The
noise-corrected share is
\begin{equation}
\omega^2_T=\frac{\mathrm{SS}_T-\mathrm{df}_T\lambda_T}{\mathrm{SS}_{\text{tot}}} .
\label{eq:omega}
\end{equation}

\textbf{Why the correction is necessary.} Without it, a factor that explains
nothing still receives a share proportional to its degrees of freedom, and
whole-plot terms receive a large one. In simulation, an uncorrected null factor
recorded a $14.6\%$ share where the truth was zero; corrected, it averaged
$-0.6\%$ over replicates. Negative values are reported as computed and never
clamped: a systematically negative share is a diagnostic that
$\hat\lambda$ is too large, and hiding it would hide the diagnostic.

\textbf{Estimating $\sigma^2_u$ requires a bias correction.} The naive pooled
estimator over $J$ region draws in $K$ groups is biased low by the factor
$c=K(J-1)/(JK-1)$ --- $0.758$ at $J=4$ --- because the pooled mean absorbs part
of the between-draw signal. Dividing by $c$ restores unbiasedness; verified over
120 replicates, recovering $0.0438$ against a true $0.0400$ where the naive
estimator gave $0.0332$.

\subsection{What each ``sense'' asks}\label{sec:senses}

Three different questions are routinely conflated under ``which variables
matter''. They are kept separate throughout.

\begin{description}
\item[Sense V --- variance-explaining.] Of the variation in a metric across
cells, how much does each design factor account for, after the noise
correction of \eqref{eq:omega}? This bounds every claim below it: a noise floor
of $50\%$ means no decomposition can explain more than half.
\item[Sense D --- decision-relevant.] Does changing a factor change
\emph{which method you would choose}? Measured as the rate at which the
best-performing method flips between adjacent factor levels, split into a raw
rate and a gated rate that counts only flips whose margin survives a
cross-fitted paired test at $z=2$. The gap between them is the quantity of
interest: it is how much of the apparent structure is noise.
\item[Sense G --- diagnostic.] For a specific method, which factors hurt
\emph{it} relative to what was achievable? Measured by
$\kappa=(\text{method}-\text{baseline})/(\text{ceiling}-\text{baseline})$, the
fraction of available headroom captured, with SuSiE as baseline and
PolyFun-oracle as ceiling.
\end{description}

Sense V and Sense D can disagree entirely, and in this dataset they do. A factor
can move a metric a great deal while never changing the winner (every method
degrades together), and a factor with a small main effect can flip the winner
if it acts differently on different methods. Reporting only one of them
overstates what is known.



\section{Completeness}\label{sec:completeness}

Every one of the 18 analysed methods has all $5{,}625$ cells populated, with
per-method declared-failure rates of $0.0\%$ and silent all-NA rates of $0.0\%$
across both arms. Three absences are structural and are \emph{not} defects:

\begin{itemize}
\item \textbf{Funmap} requires annotations and cannot run on either unannotated
arm. Its row there is \emph{empty}, not \emph{poor}; no performance claim may be
inferred from it.
\item \textbf{FINEMAP-inf} returns an all-NA posterior without raising an error
at $\phi=0.6$, $S=1$ in the sparse arm --- 45 cells, nine rows $\times$ five
region sizes. Recorded via a dedicated \texttt{bad\_pip} flag so the loss is
countable rather than averaged away.
\item \textbf{CARMA} was dropped mid-iteration on measured cost --- $19{,}000$\,s
per scenario against $9{,}765$\,s for the other eighteen combined, $66\%$ of
total runtime --- for a differentiator (LD-discrepancy outlier detection) that
an all-in-sample design cannot exercise. It is excluded from every analysis here.
\end{itemize}

\paragraph{A prior failure, and what fixed it.} An earlier execution of the
sparse-plus-infinitesimal arm recorded $38$--$50\%$ failed fits for seven
Python-backed methods. The cause was environmental, not statistical: the
toolchain lived on a filesystem subject to a rolling 30-day purge and decayed
mid-run, so activation kept succeeding while imports inside it stopped
resolving. Those results were discarded and the arm refitted after moving the
toolchain to permanent storage. The run reported here carries a per-task
preflight that imports every module the tool sources actually use --- derived by
AST from the sources rather than listed by hand --- and aborts the task if any
fails. All $1{,}800$ tasks reported \texttt{preflight ok}.

\section{Overall accuracy}\label{sec:overall}

\begin{figure}[htbp]\centering
\includegraphics[width=\linewidth]{figs/f01_ap_ranking.pdf}
\caption{Mean AP with 95\% intervals over the 5{,}625 cells of each method,
ordered by the sparse+infinitesimal continuous stratum.}\label{fig:ap}\end{figure}

Reading down the two rows of Figure~\ref{fig:ap}, three things change between
architectures.

\textbf{Everything gets harder.} Mean AP falls from $0.677$ (sparse) to $0.527$
(sparse + infinitesimal). This is expected and is not solely a difficulty
effect: the infinitesimal background is real signal that the metric counts as
noise (\S\ref{sec:design}).

\textbf{The leaderboard reorders.} Under the sparse model the PolyFun family
leads among deployable methods. Under misspecification it collapses and
\texttt{fb\_xregion} becomes the strongest non-oracle method on continuous
annotations. This is quantified in \S\ref{sec:misspec}.

\textbf{FINEMAP-inf comes into its own.} It ranks first or second across all
three sparse+infinitesimal strata. That is the architecture it explicitly
models, and the benchmark rewards it accordingly --- a useful internal check
that the second arm is doing what it was designed to do.

\section{When does each method do well?}\label{sec:conditional}

\begin{figure}[htbp]\centering
\includegraphics[width=\linewidth]{figs/f02_ap_vs_S.pdf}
\caption{AP against the number of causal variants.}\label{fig:vs-s}\end{figure}

\begin{figure}[htbp]\centering
\includegraphics[width=\linewidth]{figs/f03_ap_vs_phi.pdf}
\caption{AP against region heritability.}\label{fig:vs-phi}\end{figure}

\begin{figure}[htbp]\centering
\includegraphics[width=\linewidth]{figs/f04_ap_vs_p.pdf}
\caption{AP against nominal region size. Nearly flat in both arms --- over the
range actually swept (median 499 to 1590 variants), region size is not what
makes fine-mapping hard here.}\label{fig:vs-p}\end{figure}

\begin{figure}[htbp]\centering
\includegraphics[width=\linewidth]{figs/f05_ap_vs_pcausal.pdf}
\caption{AP against $p_{\text{causal}}$, the fraction of heritability in the
sparse component. Sparse+infinitesimal arm only.}\label{fig:vs-pc}\end{figure}

$S$ is the binding constraint in both arms: pooled AP falls
$0.780\to0.367$ from $S=1$ to $S=10$, a larger swing than any other factor
produces. $p_{\text{causal}}$ is the second lever in the infinitesimal arm,
rising $0.370\to0.640$ from $p_{\text{causal}}=0.2$ to $0.8$ --- unsurprising,
since it directly controls how much of the signal is the kind the metric scores.

\paragraph{Heritability is not monotone, and that is a real result.} Pooled AP
rises $0.482\to0.596$ from $\phi=0.05$ to $\phi=0.4$ and then \emph{falls} to
$0.586$ at $\phi=0.6$. The dip is not an artefact of the FINEMAP-inf gap --- it
survives excluding that method ($0.5907\to0.5811$) --- and it is
method-specific. Eight methods degrade at the top of the range, led by
\texttt{sbayesrc} ($-0.088$), \texttt{sparsepro} ($-0.078$) and
\texttt{finimom} ($-0.055$); ten improve, led by \texttt{susie\_inf}
($+0.032$). Methods whose priors or variational approximations are tuned for
moderate signal do not necessarily benefit from more of it. A benchmark that
reported only the pooled mean would have shown a flat top and missed this
entirely.

\section{Robustness to misspecification}\label{sec:misspec}

\begin{figure}[htbp]\centering
\includegraphics[width=\linewidth]{figs/f14_misspecification.pdf}
\caption{AP under the sparse model (blue) against sparse + infinitesimal
(orange), annotated strata. Bar length is the cost of misspecification; methods
are ordered by it.}\label{fig:misspec}\end{figure}

This is the central result of the second arm. Taking continuous annotations:

\begin{center}\small
\begin{tabular}{lrrl}\toprule
method & sparse & sparse + inf & \\\midrule
\texttt{polyfun\_est}   & \textbf{0.8061} (2nd) & 0.5480 & FDR@.9 $0.005\to0.35$\\
\texttt{polyfun\_ldsc}  & 0.7914 & 0.5443 & FDR@.9 $0.005\to0.38$\\
\texttt{fb\_xregion}    & 0.7933 (3rd) & \textbf{0.6521} (2nd) & FDR@.9 $0.049\to0.068$\\
\texttt{susie\_inf}     & 0.7398 & 0.6120 & FDR@.9 $0.005\to0.037$\\
\bottomrule\end{tabular}\end{center}

\textbf{The PolyFun family's FDR at $\pip\ge0.9$ degrades roughly seventy-fold},
from about $0.005$ to $0.35$--$0.38$, and its ranking collapses.
\texttt{fb\_xregion} loses accuracy too --- everything does --- but its FDR
moves only from $0.049$ to $0.068$, and it \emph{gains} rank, finishing second
of eighteen behind only the oracle.

The mechanism is visible in the design. PolyFun-style methods place a per-SNP
prior derived from a model in which annotation-weighted variants are causal and
everything else is null. When four-fifths of the heritability sits in unlabelled
infinitesimal effects that assumption is badly wrong, and a confident prior
built on it produces confident errors. The BEATRICE family's concrete
variational posterior does not commit as hard, and the shared prior head in
\texttt{fb\_xregion} is fitted across ten regions, which appears to regularise
it further.

\paragraph{The oracle is itself miscalibrated in this arm.} PolyFun-oracle
retains the highest AP but its FDR@$0.9$ reaches $0.36$--$0.38$, because it
holds the true \emph{sparse} prior while the infinitesimal component it does not
model generates real signal at unlabelled variants. The ceiling is a ceiling for
ranking, not for calibration, and $\kappa$ (\S\ref{sec:senseg}) inherits that
caveat.

\section{False discovery and calibration}\label{sec:calibration}

\begin{figure}[htbp]\centering
\includegraphics[width=\linewidth]{figs/f06_calibration.pdf}
\caption{Mass claimed against errors made. The whole field shifts right and up
under the infinitesimal model.}\label{fig:calib}\end{figure}

\begin{figure}[htbp]\centering
\includegraphics[width=0.82\linewidth]{figs/f07_reliability.pdf}
\caption{Reliability diagrams, pooled within arm.}\label{fig:rel}\end{figure}

\begin{figure}[htbp]\centering
\includegraphics[width=\linewidth]{figs/f08_fdr.pdf}
\caption{FDR across PIP thresholds. A calibrated method satisfies
$\mathrm{FDR}_t\le1-t$.}\label{fig:fdr}\end{figure}

In the sparse arm the trade-off is the familiar one: the BEATRICE family buys
ranking accuracy with roughly $2\times$ mass inflation and an FDR@$0.9$ an order
of magnitude above PolyFun's, while \texttt{susie\_inf} and \texttt{finemap}
stay near mass ratio $1.1$--$1.2$.

Under the infinitesimal model \emph{every} method's mass ratio rises, to
$2.5$--$4.2$. Part of this is genuine overconfidence and part is definitional:
mass placed on infinitesimal variants is real posterior belief about real
effects, counted here as excess. The two cannot be separated with this ground
truth, and the reliability diagrams in Figure~\ref{fig:rel} should be read with
that in mind --- the apparent overconfidence of the whole field in the right
panel is partly an artefact of scoring against the sparse component alone.

\textbf{What survives that caveat is the comparison between methods within the
arm}, since all are scored against the same truth. There,
\texttt{beatrice} ($2.51$) and \texttt{fb\_pooled} ($3.11$--$3.27$) are the most
conservative of the family and the PolyFun/SuSiE group the least
($4.2$).

\begin{figure}[htbp]\centering
\includegraphics[width=\linewidth]{figs/f17_murphy.pdf}
\caption{Murphy decomposition. REL is miscalibration (lower better); RES is
resolution (higher better). RES is what distinguishes an informative posterior
from one calibrated by being vague.}\label{fig:murphy}\end{figure}

\section{Credible sets}\label{sec:cs}

\begin{figure}[htbp]\centering
\includegraphics[width=\linewidth]{figs/f09_credible_sets.pdf}
\caption{95\% credible sets: follow-up burden against coverage.}\label{fig:cs}\end{figure}

Size and coverage are always read together: a set can be made to cover by being
large and small by abandoning coverage. Neither number alone identifies which
situation a method is in.
\section{Sense V: what moves the metric}\label{sec:sensev}

\subsection{How the decomposition is computed, in full}\label{sec:vhow}

Every number in this section descends from the per-fit AP values of
\S\ref{sec:ap} by the following chain. Constants are named as they appear in the
code: \texttt{N\_ITER}$=10$, \texttt{N\_REGIONS}$=2$ (regions per size class),
\texttt{N\_FITS}$=20$, \texttt{M\_SHARED}$=25$, \texttt{SUBPLOT\_FACTORS}$=\{S,\phi\}$.

\paragraph{Step 1 --- form cell means.}
A cell is the tuple $(\text{job\_dir},S,\phi,\text{region size})$. Its value is
the mean of the $\texttt{N\_FITS}=20$ fits it contains: 10 iterations $\times$ 2
regions of that size class. This table (L3) is the only input to the ANOVA;
individual fits are used solely to estimate the two variance components below.

\paragraph{Step 2 --- estimate $\sigma^2_\varepsilon$, the iteration variance.}
Group the fit table by
\[
\mathrm{KEY}=(\text{job\_dir},\,\text{region size},\,\text{region\_idx},\,S,\,\phi),
\]
which pins down \emph{one physical region at one design point}, leaving only the
iteration index free. Within each group take the sample variance across its 10
iterations; $\hat\sigma^2_\varepsilon$ is the mean of those variances. Note what
is held fixed: the region, its LD, its variants and its annotations are
literally identical across those 10 values, so this is pure replicate error from
the phenotype draw.

\paragraph{Step 3 --- estimate $\sigma^2_u$, the between-region-draw variance.}
This cannot be read off directly, because a region draw is never replicated
within a size class --- there are exactly two, and they differ. It is recovered
from their contrast:

\begin{enumerate}
\item Average over iterations within each KEY group, giving a region mean
$\bar y(\text{job},\text{size},\text{idx},S,\phi)$.
\item Reshape wide on \texttt{region\_idx} and difference the pair:
$d = \bar y_1 - \bar y_2$ for each $(\text{job},\text{size},S,\phi)$.
\item In expectation
$\operatorname{Var}(d) = 2\sigma^2_u + 2\sigma^2_\varepsilon/\texttt{N\_ITER}$,
so subtracting the second term and halving would give $\sigma^2_u$ --- except
for the bias in Step 4.
\item Do this separately for each region size, then average the five estimates
into $\bar\sigma^2_u$.
\end{enumerate}

\paragraph{Step 4 --- debias the pair variance.}
The two region draws of a size class are \emph{reused} across all
$K$ combinations of $(S,\phi)$ within a job, so those $K$ differences share a
common draw-level offset $D_j$ and are not independent. With $J$ job rows and
$K$ cells per job the naive $\operatorname{Var}(d)$ is biased low by
\[
c=\frac{K(J-1)}{JK-1},
\qquad
\hat\sigma^2_D=\max\!\Bigl(0,\ \frac{\operatorname{Var}(d)-2\sigma^2_\varepsilon/\texttt{N\_ITER}}{c}\Bigr),
\qquad
\hat\sigma^2_u=\hat\sigma^2_D/2 .
\]
At $J=4$ this is $c=0.758$, i.e.\ the naive estimator recovers only about
three-quarters of the true component. Verified over 120 simulated replicates:
corrected $0.0438$ against a true $0.0400$, where the naive estimator gave
$0.0332$.

\paragraph{Step 5 --- fit the saturated ANOVA on cell means.}
\texttt{aov(response \textasciitilde{} f1 * f2 * \ldots)} over the stratum's free
factors, all interactions included. The design is complete and balanced, so the
model is saturated: \textbf{there is deliberately no residual line}, every degree
of freedom belongs to a named term, and the orthogonal sums of squares partition
$\mathrm{SS}_{\text{tot}}$ exactly. If a \texttt{Residuals} row does appear the
code warns and drops it, because its presence means the stratum is not balanced
and the closed forms below no longer hold.

\paragraph{Step 6 --- classify each term.}
\[
T \text{ is \emph{whole-plot}} \iff T \text{ contains \textbf{no} factor from } \{S,\phi\}.
\]
The logic runs through what a contrast differences away. $S$ and $\phi$ are
re-simulated per scenario \emph{on the same region draw}, so any contrast
involving them cancels the draw exactly and sees only iteration noise. A term
built solely from region size and enrichment compares values computed on
\emph{different} region draws and therefore carries $\sigma^2_u$ as well. Note
the consequence: \texttt{enrich} alone and
\texttt{enrich:region\_size} are whole-plot, whereas \texttt{S:region\_size} is
not, despite containing region size.

\paragraph{Step 7 --- the noise expectation and the corrected share.}
\[
\lambda_T=\frac{\sigma^2_\varepsilon}{\texttt{N\_FITS}}
+\mathbf{1}\{T\text{ whole-plot}\}\cdot\frac{\texttt{M\_SHARED}\cdot\bar\sigma^2_u}{2},
\qquad
\omega^2_T=\frac{\mathrm{SS}_T-\mathrm{df}_T\lambda_T}{\mathrm{SS}_{\text{tot}}} .
\]
\texttt{N\_FITS}$=20$ is the number of fits behind each cell mean;
\texttt{M\_SHARED}$=25=5\times5$ is the number of $(S,\phi)$ cells sharing one
region draw, and the $/2$ is the two regions per size class. The noise floor is
$\sum_T \mathrm{df}_T\lambda_T/\mathrm{SS}_{\text{tot}}$, and it is reported
split into its iteration and region parts.

\paragraph{Why cell means rather than one replicate-level ANOVA.}
Both are defensible in principle; only one is safe here. At replicate level each
noise source enters a mean square multiplied by the number of rows sharing that
draw --- hundreds of them --- so subtracting a cell-mean-scale $\sigma^2_u/2$
under-corrects by orders of magnitude \emph{while still looking like a
correction has been applied}. Working on cell means keeps every quantity on one
scale, which is why the closed form above is exact rather than approximate.

\paragraph{Sobol indices.}
$V_T=\max(\omega^2_T,0)$ --- this is the \emph{only} place a clamp is applied,
and it is applied because the indices are ratios that would otherwise be
uninterpretable. Then
$S_i = V_i/\sum_T V_T$ and $S_{Ti}=\sum_{T\ni i}V_T/\sum_T V_T$, so
$S_{Ti}-S_i$ is the share of variance factor $i$ participates in only through
interactions. Everywhere else in this report negative $\omega^2$ is reported as
computed.


\subsection{The variance budget}

\begin{figure}[htbp]\centering
\includegraphics[width=\linewidth]{figs/f10_variance_budget.pdf}
\caption{The closed variance budget for AP. Strata whose design contains a
single row have no identifiable $\sigma^2_u$ and are omitted.}\label{fig:budget}\end{figure}

The noise floor runs $17$--$19\%$ in the strata where it is estimable, with main
effects at $60$--$68\%$ and everything above two-way under $5\%$. That is the
licence for reading the marginal plots of \S\ref{sec:conditional} at face value.

\paragraph{Two strata have no budget at all.} \texttt{sparse/none} contains a
single design row, so there is only one group from which to estimate
$\sigma^2_u$ and the between-region-draw component is unidentifiable. Every
corrected share in this section therefore excludes it. Uncorrected shares would
be inflated, particularly for whole-plot terms, and are not reported.

\subsection{Different metrics have different drivers}

\begin{figure}[htbp]\centering
\includegraphics[width=\linewidth]{figs/f11_sense_v_responses.pdf}
\caption{Noise-corrected $\omega^2$ for three responses. The factor ordering
changes with the response.}\label{fig:vresp}\end{figure}

Ranking accuracy and mass calibration are both governed by $S$; false discovery
is governed by $\phi$. A benchmark reporting AP alone would call $\phi$
second-order, when for the metric a practitioner acts on --- the error rate
among confident calls --- it is the leading factor.

\paragraph{Negative shares are reported unclamped.} Where
$\mathrm{df}_T\lambda_T$ exceeds $\mathrm{SS}_T$ the corrected share goes
negative, which says the estimated noise for that term is larger than the
variation observed. FDR at a fixed threshold is a ratio of small counts and is
far noisier per cell than AP, so its decomposition is unreliable at this
replication level. Clamping would hide exactly that diagnostic.

\section{Sense D: what changes the decision}\label{sec:sensed}

\subsection{How a cell's winner is decided, in full}\label{sec:dhow}

The unit is again the cell $(\text{job\_dir},S,\phi,\text{region size})$, but
here it is \emph{not} collapsed to a mean. The procedure per cell is:

\paragraph{Step 1 --- build the replicate matrix.}
Reshape to a matrix $Y$ with one row per iteration ($R=10$) and one column per
method, entry $Y_{im}=\mathrm{AP}$ of method $m$ on iteration $i$. Because all
methods ran on the same simulated data, row $i$ is a genuine paired
observation across methods. Columns that are entirely non-finite are dropped;
fewer than two surviving columns means the cell is skipped.

\paragraph{Step 2 --- the nominal winner.}
$\text{winner}=\arg\max_m \overline{Y}_{\cdot m}$, the column mean over all 10
iterations. This is the quantity the regime maps display, and on its own it is
\textbf{not} evidence of anything --- the maximum of $M$ noisy means is upward
biased by construction.

\paragraph{Step 3 --- cross-fitted testing.}
If $R<4$ the cell is marked undecided immediately. Otherwise the 10 iterations
are randomly permuted (once, under \texttt{set.seed(20260806)}, so the result is
reproducible) and split into halves $A$ and $B$. Then, in each direction:
\begin{enumerate}
\item Compute column means on the \emph{selection} half.
\item Take the top two methods \emph{by that half only}.
\item On the \emph{other} half form the paired differences
$d_i = Y_{i,\text{top1}} - Y_{i,\text{top2}}$.
\item Declare decided iff $|\bar d| > z\cdot\mathrm{SE}(\bar d)$ with $z=2$ and
$\mathrm{SE}=\mathrm{sd}(d)/\sqrt{n}$.
\end{enumerate}

\paragraph{Step 4 --- require agreement.}
The cell counts as decided only if \textbf{both} directions ($A$ selects and $B$
tests, then $B$ selects and $A$ tests) declare decided \emph{and} both name the
same winner. Any disagreement means the apparent winner depended on which half
happened to be used for selection.

\paragraph{Why cross-fitting is necessary.}
Selecting the top two and testing them on the same data guarantees a positive
margin: the selection has already sorted the noise. Splitting selection from
testing removes that bias, and requiring both directions to agree removes the
residual luck of a single split. This is why the gated rates are so much lower
than the raw ones --- the raw rate is measuring exactly the bias this procedure
is designed to eliminate.

\paragraph{Step 5 --- the reported margin.}
$\bar d$ (\texttt{dbar}) is computed on \emph{all} 10 rows as
$\overline{Y_{\cdot\text{winner}} - Y_{\cdot\text{runner-up}}}$, using the
cross-fitted runner-up. It is a descriptive effect size, not the test statistic.

\paragraph{Step 6 --- flip rates.}
For a factor $f$: hold every other free factor fixed, walk adjacent levels of
$f$, and record whether the winner changes. \texttt{flip\_raw} is the fraction
of such adjacent pairs whose nominal winner differs; \texttt{flip\_gated} is the
same fraction counting only pairs in which the relevant cells were decided.
\texttt{pct\_cells\_decided} is reported separately as the unconditional
decidability of the stratum.


\subsection{Who actually wins}

\begin{figure}[htbp]\centering
\includegraphics[width=\linewidth]{figs/f12_sense_d_winners.pdf}
\caption{Cells won by each method, split by whether the win survives the
cross-fitted paired test.}\label{fig:winners}\end{figure}

Grey dominates every panel. The nominal winner of a cell usually cannot be
separated from its runner-up, and win counts must not be read as a ranking ---
most are ties dressed up as wins.

\paragraph{What follows.} At 20 fits per cell the benchmark cannot resolve which
of eighteen methods is best \emph{within a single cell}, and regime maps built
from modal winners are not recommendations. This is a statement about power, not
about the methods being equivalent: the paired contrasts of
\S\ref{sec:senseg} pool tens of thousands of replicates and are decisive.
Reporting only one of the two would misrepresent what the benchmark establishes.

\section{Sense G: headroom and the ablation}\label{sec:senseg}

\subsection{How $\kappa$ and the controlled pairs are computed, in full}\label{sec:ghow}

The input here is the \emph{replicate-level} table (L2): one row per
$(\text{cell},\text{iteration},\text{method})$, reshaped so that each
$(\text{cell},\text{iteration})$ carries one column per method. Because every
method ran on the same simulated data, a row is a set of paired observations ---
same causal variants, same genotypes, same annotations, same phenotype draw.

\paragraph{Step 1 --- headroom.}
Per row,
\[
h=\mathrm{AP}_{\texttt{polyfun\_oracle}}-\mathrm{AP}_{\texttt{susie}} .
\]
\texttt{susie} is the annotation-blind backbone --- what a method achieves
\emph{without} using annotations at all --- and \texttt{polyfun\_oracle} receives
the simulator's true prior weights, so it is what is achievable \emph{with
perfect} annotation information. Their difference is the room annotations
create on that particular replicate.

\paragraph{Step 2 --- the gate.}
Rows with $h\le\texttt{HEADROOM\_EPS}=0.02$ (or non-finite $h$) are excluded,
because $\kappa$ divides by $h$ and is unstable as $h\to0$. The excluded
fraction and the median headroom are reported rather than dropped silently
(\S\ref{sec:headroom}); on the unannotated arm the gate removes
\emph{everything}, which is correct: the ceiling and the baseline are the same
method there, so $h\equiv0$ and $\kappa$ is undefined by construction.

\paragraph{Step 3 --- $\kappa$.}
On surviving rows, for each focus method $m\in\{$\texttt{functional\_beatrice},
\texttt{fb\_xregion}, \texttt{fb\_pooled}$\}$,
\[
\kappa_m=\frac{\mathrm{AP}_m-\mathrm{AP}_{\texttt{susie}}}{h} .
\]
$\kappa=1$ means $m$ captured all the available headroom; $\kappa=0$ means it
did no better than ignoring the annotations; $\kappa<0$ means it did worse than
the annotation-blind baseline on that replicate.

\paragraph{Step 4 --- the controlled pairs, which are NOT gated.}
Two contrasts are formed on \emph{all} rows, with no headroom gate, because they
are differences rather than ratios and have no denominator to destabilise:
\[
\Delta_{\text{annot}}=\mathrm{AP}_{\texttt{functional\_beatrice}}-\mathrm{AP}_{\texttt{beatrice}},
\qquad
\Delta_{\text{joint}}=\mathrm{AP}_{\texttt{fb\_xregion}}-\mathrm{AP}_{\texttt{functional\_beatrice}} .
\]
Each isolates exactly one mechanism, because the two methods being differenced
are identical in every other respect: same inference network, same
hyperparameters, same data, same seed. $\Delta_{\text{annot}}$ turns the
annotation prior on; $\Delta_{\text{joint}}$ changes that prior from per-region
to shared-across-regions.

\paragraph{Step 5 --- inference on the pairs.}
For a group of $n$ paired rows, $\bar\Delta$, $\mathrm{sd}(\Delta)$,
$\mathrm{SE}=\mathrm{sd}/\sqrt n$ and $t=\bar\Delta/\mathrm{SE}$. The pairing is
what supplies the power: the between-scenario variance that dominates marginal
comparisons cancels exactly, leaving only the within-replicate difference. This
is why $\S\ref{sec:senseg}$ resolves effects of $0.02$ AP decisively while
\S\ref{sec:sensed} cannot resolve a $0.13$ gap between two \emph{different}
methods in a single cell --- $n$ differs by three orders of magnitude and the
nuisance variance has been removed.

\paragraph{Step 6 --- conditional breakdowns.}
The figures in \S\ref{sec:where} are $\bar\Delta$ recomputed within levels of one
design factor at a time, pooling the two annotated strata. $114$ of the $22{,}500$
$\Delta$ values are missing ($0.5\%$), all at $S\in\{5,10\}$, corresponding to
the BEATRICE family's small fit-failure rate; those rows are dropped pairwise.


\subsection{Headroom exists in both arms}\label{sec:headroom}

\begin{center}\small
\begin{tabular}{llrrr}\toprule
arm & annotation & median headroom & \% $\le0$ & replicate rows\\\midrule
sparse & none & 0.0000 & 100.0\% & 1{,}250\\
sparse & binary & 0.0840 & 25.5\% & 5{,}000\\
sparse & continuous & 0.1662 & 20.6\% & 5{,}000\\
sparse + inf & none & 0.0000 & 100.0\% & 5{,}000\\
sparse + inf & binary & 0.0667 & 21.2\% & 20{,}000\\
sparse + inf & continuous & \textbf{0.1915} & 13.9\% & 20{,}000\\
\bottomrule\end{tabular}\end{center}

Despite the oracle's poor calibration under misspecification
(\S\ref{sec:misspec}), its \emph{ranking} advantage over the baseline is intact,
so $\kappa$ remains well defined in the annotated strata. The largest median
headroom of any stratum is sparse+infinitesimal / continuous at $0.1915$ ---
there is \emph{more} for annotations to exploit when the model is misspecified,
not less. Both unannotated arms have headroom identically zero, so the gate
removes them entirely; that is correct, since the ceiling and the baseline are
the same method there.

\subsection{The ablation}

\begin{figure}[htbp]\centering
\includegraphics[width=\linewidth]{figs/f13_ablation.pdf}
\caption{Paired contrasts with 95\% intervals, all six strata.}\label{fig:ablation}\end{figure}

\begin{center}\small
\begin{tabular}{llrrrr}\toprule
arm & annotation & annotations & region sharing & combined & $n$\\\midrule
sparse & none & $0.0000$ & $0.0000$ & $0.0000$ & 1{,}250\\
sparse & binary & $+0.0254$ & $+0.0205$ & $+0.0458$ & 4{,}980\\
sparse & continuous & $+0.0224$ & $+0.0730$ & $+0.0954$ & 4{,}960\\
sparse+inf & none & $0.0000$ & $0.0000$ & $0.0000$ & 5{,}000\\
sparse+inf & binary & $+0.0149$ & $+0.0096$ & $+0.0244$ & 20{,}000\\
sparse+inf & continuous & $+0.0116$ & $\mathbf{+0.0522}$ & $+0.0638$ & 20{,}000\\
\bottomrule\end{tabular}\end{center}

All twelve annotated contrasts are positive with $t$ between $17$ and $82$
(Table~\ref{tab:ablation}). Both unannotated arms are \textbf{exactly
$0.0000\pm0.0000$} over $6{,}250$ paired replicates --- the three variants are
bit-identical there, so the measured gains cannot be an artefact of the extra
machinery, of additional parameters, or of run-to-run variation.

Gains are smaller under the infinitesimal model, which is expected: the
background is variance no annotation can explain. But the \emph{ordering}
inverts in an informative way --- region sharing contributes four times what
annotations alone do on continuous annotations there ($+0.052$ against
$+0.012$), against roughly three times in the sparse arm. \textbf{The harder and
more misspecified the problem, the more the sharing matters relative to the
annotations themselves.}

\subsection{Where the gains come from}\label{sec:where}

\begin{figure}[htbp]\centering
\includegraphics[width=\linewidth]{figs/f15_ablation_by_S.pdf}
\caption{Both contrasts against $S$.}\label{fig:abl-s}\end{figure}

\begin{figure}[htbp]\centering
\includegraphics[width=\linewidth]{figs/f16_ablation_by_phi.pdf}
\caption{Both contrasts against $\phi$.}\label{fig:abl-phi}\end{figure}

Region sharing is worth most where regions contain fewest causals and decays as
$S$ grows. That is what the mechanism predicts: a region with a single causal
variant carries almost no internal information about which annotations matter,
so borrowing the enrichment estimate from nine other regions is worth a great
deal, while a region with ten causals can estimate it locally.

\subsection{$\kappa$ is reported but no conclusion rests on it}

After gating, $\kappa$'s standard deviations exceed unity against means well
below it. The denominator is small in many surviving cells and the ratio is
correspondingly unstable. Every conclusion in this section rests on the
controlled pairs, which are differences rather than ratios.

\paragraph{The $\kappa$ variance decomposition is unavailable.} The
replicate-level table does not carry the region index, so $\sigma^2_u$ cannot be
estimated for the derived $\kappa$ and $\Delta$ responses. The code refuses to
report an uncorrected decomposition, which would inflate whole-plot terms by
$m\sigma^2_u/2$ per degree of freedom.

\section{Validity checks}\label{sec:validity}

Mechanically verified on this run; all pass.

\begin{enumerate}
\item In-sample LD only --- zero rows carrying a reference-panel size.
\item 45 design rows present, matching the grid.
\item Balanced cells per analysed method --- exactly $5{,}625$ for all 18.
\item CARMA excluded by design and reported as such.
\item No method entirely absent outside the unannotated arms.
\item \texttt{polyfun\_est} $=$ \texttt{susie} on both unannotated arms, exactly.
\item \texttt{polyfun\_ldsc} $=$ \texttt{susie} on both unannotated arms, exactly.
\end{enumerate}

Additional checks performed for this report, beyond the pipeline's own:

\begin{enumerate}\setcounter{enumi}{7}
\item \textbf{Truth consistency}: $\texttt{n\_causal}=S$ for all $2{,}031{,}230$
fit rows; zero mismatches.
\item \textbf{Grid completeness}: exactly $11{,}250$ distinct
$(\text{row},S,\phi,\text{iter})$ combinations; \texttt{region\_idx} takes only
the values 1 and 2.
\item \textbf{Metric internal consistency}: $\mathrm{tp}_t\le n_t$ and
$\mathrm{tp}_t\le S$ at every threshold; $n_t$ monotone decreasing in $t$;
$\mathrm{AP}\in[0,1]$; $|CS_{95}|\le p$; band counts summing to $p$; causal band
counts summing to $S$. Zero violations in every case.
\item \textbf{Negative controls, exactly}: all four BEATRICE variants identical
to machine precision on both unannotated arms (max $|{\rm diff}|=0$ over 625
cells), and \texttt{polyfun\_\{est,ldsc,oracle\}} identical to \texttt{susie}
there.
\item \textbf{Monotonicity where it should hold}: AP strictly decreasing in $S$;
strictly increasing in $p_{\text{causal}}$; sparse arm easier than
sparse+infinitesimal. AP is \emph{not} monotone in $\phi$, which
\S\ref{sec:conditional} shows to be a genuine method-specific effect rather than
a fault.
\item \textbf{Simulation code}: PVE calibration verified to be exact by
construction in both architectures; the oracle prior verified to use the
identical formula to the causal-selection weights.
\end{enumerate}

\section{Limitations}\label{sec:limits}

\begin{enumerate}
\item \textbf{In-sample LD only.} Iteration 003 established that in-sample and
reference-panel LD can invert method rankings. Nothing here transfers to the
reference-panel setting.
\item \textbf{The infinitesimal arm scores the sparse component only.} Non-causal
variants carry real effects that count as false discoveries. Absolute FDR is not
comparable across arms; degradation between arms is.
\item \textbf{Nominal region sizes are approximate above 1000}, with the 2000
class median 1590.
\item \textbf{Range dependence.} All importance statements are relative to the
ranges swept.
\item \textbf{Per-cell decisions are underpowered} (\S\ref{sec:sensed}).
\item \textbf{Two strata lack a noise floor} and are excluded from Sense V.
\item \textbf{$\kappa$ decomposition unavailable}; the replicate table lacks the
region index.
\item \textbf{One tuning configuration.} BEATRICE-family hyperparameters come
from a single Optuna study and are held fixed; competitors run at documented
defaults. A separate analysis found the family's sensitivity to those
hyperparameters to be small ($\pm0.004$ AP across a wide hand-tuned sweep)
relative to the architectural contrasts reported here.
\end{enumerate}

\section{Full results tables}\label{sec:tables}

Bold rows are the BEATRICE family. \emph{cells} counts non-missing cells.

\begin{table}[htbp]\centering\small
\caption{\textbf{No annotations} --- all methods, sparse arm. Ranked by mean AP. \emph{cells} is the number of non-missing cells out of 1{,}125; anything below that is a genuine absence, explained in \S\ref{sec:completeness}.}
\label{tab:main-none}
\begin{tabular}{lrrrrrrrr r}\toprule
method & AP & FDR@.9 & mass & ECE$_{hi}$ & $|CS_{95}|$ & cov$_{95}$ & RES & BSS & cells\\\midrule
susie\_inf & 0.7095 & 0.0062 & 1.14 & 0.0288 & 4.2 & 19.864 & 0.0023 & 0.575 & 125\\
polyfun\_est & 0.6819 & 0.0061 & 1.73 & 0.0225 & 1.7 & 19.800 & 0.0024 & 0.579 & 125\\
polyfun\_ldsc & 0.6819 & 0.0061 & 1.73 & 0.0225 & 1.7 & 19.800 & 0.0024 & 0.579 & 125\\
polyfun\_oracle & 0.6819 & 0.0061 & 1.73 & 0.0225 & 1.7 & 19.800 & 0.0024 & 0.579 & 125\\
susie & 0.6819 & 0.0061 & 1.73 & 0.0225 & 1.7 & 19.800 & 0.0024 & 0.579 & 125\\
finemap & 0.6767 & 0.0442 & 1.16 & 0.0683 & 1.5 & 19.168 & 0.0021 & 0.490 & 125\\
finemap\_inf & 0.6711 & 0.1815 & 1.82 & 0.1619 & 1.7 & 16.864 & 0.0021 & 0.056 & 120\\
\textbf{beatrice} & 0.6659 & 0.0237 & 1.84 & 0.1434 & 1.9 & 18.888 & 0.0018 & 0.321 & 125\\
\textbf{fb\_pooled} & 0.6659 & 0.0237 & 1.84 & 0.1434 & 1.9 & 18.888 & 0.0018 & 0.321 & 125\\
\textbf{fb\_xregion} & 0.6659 & 0.0237 & 1.84 & 0.1434 & 1.9 & 18.888 & 0.0018 & 0.321 & 125\\
\textbf{functional\_beatrice} & 0.6659 & 0.0237 & 1.84 & 0.1434 & 1.9 & 18.888 & 0.0018 & 0.321 & 125\\
sparsepro & 0.5737 & 0.1556 & 1.45 & 0.1232 & 1.5 & 17.336 & 0.0019 & 0.249 & 125\\
finimom & 0.5735 & 0.1213 & 1.12 & 0.0890 & 1.7 & 18.056 & 0.0019 & 0.322 & 125\\
paintor & 0.5247 & 0.0570 & 0.74 & 0.0615 & 1.8 & 19.008 & 0.0013 & 0.372 & 125\\
sbayesrc & 0.4841 & 0.3705 & 10.02 & 0.2207 & 1.0 & 13.064 & 0.0015 & -0.552 & 125\\
marginal\_z & 0.4086 & 0.0000 & 0.43 & NaN & 735.6 & 20.000 & 0.0000 & -0.003 & 125\\
abf & 0.4005 & 0.0054 & 0.43 & 0.0237 & 3.3 & 19.920 & 0.0008 & 0.314 & 125\\
funmap & NaN & 0.0000 & NaN & NaN & NaN & 0.000 & NaN & NaN & 0\\
\bottomrule\end{tabular}\end{table}
\begin{table}[htbp]\centering\small
\caption{\textbf{Binary annotations} --- all methods, sparse arm. Ranked by mean AP. \emph{cells} is the number of non-missing cells out of 1{,}125; anything below that is a genuine absence, explained in \S\ref{sec:completeness}.}
\label{tab:main-binary}
\begin{tabular}{lrrrrrrrr r}\toprule
method & AP & FDR@.9 & mass & ECE$_{hi}$ & $|CS_{95}|$ & cov$_{95}$ & RES & BSS & cells\\\midrule
polyfun\_oracle & 0.8432 & 0.0110 & 1.63 & 0.0715 & 1.2 & 19.866 & 0.0032 & 0.694 & 500\\
polyfun\_est & 0.7914 & 0.0046 & 1.81 & 0.0376 & 1.4 & 19.914 & 0.0028 & 0.644 & 500\\
polyfun\_ldsc & 0.7732 & 0.0209 & 1.81 & 0.0440 & 1.3 & 19.668 & 0.0027 & 0.635 & 500\\
\textbf{fb\_xregion} & 0.7484 & 0.0569 & 2.08 & 0.1454 & 1.3 & 19.306 & 0.0023 & 0.301 & 500\\
susie\_inf & 0.7470 & 0.0043 & 1.16 & 0.0283 & 2.1 & 19.870 & 0.0026 & 0.620 & 500\\
funmap & 0.7462 & 0.0866 & 1.81 & 0.0746 & 1.3 & 18.616 & 0.0026 & 0.614 & 500\\
\textbf{functional\_beatrice} & 0.7280 & 0.0783 & 2.10 & 0.1547 & 1.3 & 19.034 & 0.0022 & 0.245 & 500\\
susie & 0.7232 & 0.0050 & 1.80 & 0.0218 & 1.4 & 19.856 & 0.0026 & 0.624 & 500\\
finemap & 0.7147 & 0.0447 & 1.14 & 0.0648 & 1.3 & 19.346 & 0.0023 & 0.535 & 500\\
finemap\_inf & 0.7096 & 0.1717 & 1.76 & 0.1585 & 1.5 & 16.894 & 0.0024 & 0.141 & 480\\
\textbf{beatrice} & 0.7026 & 0.0235 & 1.79 & 0.1592 & 1.9 & 19.050 & 0.0021 & 0.369 & 500\\
\textbf{fb\_pooled} & 0.6661 & 0.0745 & 2.09 & 0.1547 & 1.5 & 18.312 & 0.0019 & 0.187 & 500\\
paintor & 0.6458 & 0.0668 & 0.74 & 0.0919 & 1.1 & 19.064 & 0.0016 & 0.452 & 500\\
finimom & 0.6182 & 0.1156 & 1.09 & 0.0876 & 1.4 & 18.264 & 0.0022 & 0.378 & 500\\
sparsepro & 0.6107 & 0.1509 & 1.45 & 0.1232 & 1.3 & 17.504 & 0.0021 & 0.286 & 500\\
sbayesrc & 0.5086 & 0.3592 & 4.74 & 0.2470 & 1.0 & 13.186 & 0.0018 & -0.313 & 500\\
marginal\_z & 0.4467 & 0.0000 & 0.43 & NaN & 765.5 & 20.000 & 0.0000 & -0.003 & 500\\
abf & 0.4357 & 0.0047 & 0.43 & 0.0255 & 2.7 & 19.914 & 0.0009 & 0.334 & 500\\
\bottomrule\end{tabular}\end{table}
\begin{table}[htbp]\centering\small
\caption{\textbf{Continuous annotations} --- all methods, sparse arm. Ranked by mean AP. \emph{cells} is the number of non-missing cells out of 1{,}125; anything below that is a genuine absence, explained in \S\ref{sec:completeness}.}
\label{tab:main-continuous}
\begin{tabular}{lrrrrrrrr r}\toprule
method & AP & FDR@.9 & mass & ECE$_{hi}$ & $|CS_{95}|$ & cov$_{95}$ & RES & BSS & cells\\\midrule
polyfun\_oracle & 0.9114 & 0.0119 & 1.21 & 0.0791 & 1.0 & 19.952 & 0.0037 & 0.774 & 500\\
polyfun\_est & 0.8061 & 0.0051 & 1.77 & 0.0414 & 1.6 & 19.928 & 0.0028 & 0.641 & 500\\
\textbf{fb\_xregion} & 0.7933 & 0.0491 & 2.14 & 0.1529 & 1.3 & 19.462 & 0.0027 & 0.327 & 500\\
polyfun\_ldsc & 0.7914 & 0.0047 & 1.77 & 0.0360 & 1.6 & 19.906 & 0.0028 & 0.638 & 500\\
\textbf{fb\_pooled} & 0.7681 & 0.0422 & 2.15 & 0.1504 & 1.6 & 19.424 & 0.0023 & 0.278 & 500\\
funmap & 0.7617 & 0.0227 & 1.77 & 0.0367 & 1.6 & 19.610 & 0.0027 & 0.628 & 500\\
susie\_inf & 0.7398 & 0.0053 & 1.13 & 0.0284 & 3.0 & 19.886 & 0.0026 & 0.618 & 500\\
\textbf{functional\_beatrice} & 0.7204 & 0.1002 & 2.27 & 0.1728 & 1.2 & 18.680 & 0.0022 & 0.117 & 500\\
susie & 0.7175 & 0.0059 & 1.77 & 0.0236 & 1.6 & 19.852 & 0.0027 & 0.621 & 500\\
finemap & 0.7071 & 0.0492 & 1.16 & 0.0682 & 1.4 & 19.250 & 0.0023 & 0.529 & 500\\
finemap\_inf & 0.7053 & 0.1708 & 1.77 & 0.1555 & 1.6 & 16.870 & 0.0024 & 0.143 & 480\\
paintor & 0.7001 & 0.0520 & 0.74 & 0.0790 & 1.0 & 19.420 & 0.0019 & 0.520 & 500\\
\textbf{beatrice} & 0.6981 & 0.0198 & 1.82 & 0.1531 & 1.9 & 18.846 & 0.0021 & 0.362 & 500\\
sbayesrc & 0.6192 & 0.3000 & 5.75 & 0.1993 & 1.0 & 14.482 & 0.0024 & -0.254 & 500\\
finimom & 0.6177 & 0.1152 & 1.09 & 0.0910 & 1.6 & 18.308 & 0.0022 & 0.378 & 500\\
sparsepro & 0.6093 & 0.1502 & 1.45 & 0.1236 & 1.4 & 17.636 & 0.0021 & 0.294 & 500\\
marginal\_z & 0.4371 & 0.0000 & 0.43 & NaN & 738.4 & 20.000 & 0.0000 & -0.003 & 500\\
abf & 0.4271 & 0.0024 & 0.43 & 0.0237 & 3.0 & 19.910 & 0.0009 & 0.333 & 500\\
\bottomrule\end{tabular}\end{table}
\begin{table}[htbp]\centering\small
\caption{Paired ablation on AP. Each row compares two methods on the \emph{same} simulated replicate; $n$ is the number of paired replicates. $t=\bar d/\mathrm{SE}(\bar d)$.}\label{tab:ablation}
\begin{tabular}{llrrrrr}\toprule
stratum & contrast & $\bar d$ & sd & $n$ & SE & $t$\\\midrule
No annotations & annotations & 0.0000 & 0.0000 & 1250 & 0.00000 & ---\\
No annotations & region sharing & 0.0000 & 0.0000 & 1250 & 0.00000 & ---\\
No annotations & combined & 0.0000 & 0.0000 & 1250 & 0.00000 & ---\\
Binary annotations & annotations & 0.0254 & 0.0975 & 4981 & 0.00138 & 18.4\\
Binary annotations & region sharing & 0.0205 & 0.0859 & 4980 & 0.00122 & 16.8\\
Binary annotations & combined & 0.0458 & 0.1026 & 4980 & 0.00145 & 31.5\\
Continuous annotations & annotations & 0.0224 & 0.0962 & 4965 & 0.00136 & 16.4\\
Continuous annotations & region sharing & 0.0730 & 0.1183 & 4960 & 0.00168 & 43.5\\
Continuous annotations & combined & 0.0954 & 0.1352 & 4960 & 0.00192 & 49.7\\
\bottomrule\end{tabular}\end{table}


\end{document}
